%!TEX program = xelatex
\PassOptionsToPackage{full}{textcomp}
\documentclass[lang=en,11pt,a4paper,citestyle=authoryear]{elegantpaper}

\title{HW4 - STAT 709}
\author{Wells Guan}

\usepackage{array,url,stix}
\usepackage{tikz}
\usepackage{tikz-cd}
\usepackage{enumitem}
\usepackage{amsmath}
\usepackage{amssymb}
\usepackage{listings}
\usepackage{subfigure}

\lstset{
    frame=single,
    columns=fullflexible,
    keepspaces=true,
    showstringspaces=false,
    breaklines=true,
    breakatwhitespace=false,
    basicstyle=\ttfamily\footnotesize,
    numbers=left,
    numberstyle=\tiny,
    numbersep=6pt
}

\newcommand{\Z}{\mathbb{Z}}
\newcommand{\R}{\mathbb{R}}
\newcommand{\N}{\mathbb{N}}
\newcommand{\B}{\mathcal{B}}
\newcommand{\E}{\mathbb{E}}
\renewcommand{\P}{\mathbb{P}}
\newcommand{\norm}[1]{\left\lVert#1\right\rVert}

\begin{document}
\maketitle

% The assignment PDF gives only the textbook exercise numbers.
% Insert the exact textbook statements here if needed.
\subsection*{Problem 1. Textbook exercises}
Jun Shao's \textit{Mathematical Statistics}, Chapter 2:
Exercise 31 (required); Exercise 25 (optional).

\vspace{1em}

\subsubsection*{Ex. 2.31}
Let $X_1,\ldots,X_n$ be i.i.d.\ random variables having a distribution
$P\in\mathcal{P}$, where $\mathcal{P}$ is the family of distributions on
$\mathbb{R}$ having continuous c.d.f.'s. Let
\[
T=\bigl(X_{(1)},\ldots,X_{(n)}\bigr)
\]
be the vector of order statistics. Show that, given $T$, the conditional
distribution of
\[
X=(X_1,\ldots,X_n)
\]
is a discrete distribution putting probability $1/n!$ on each of the
$n!$ points
\[
\bigl(X_{(i_1)},\ldots,X_{(i_n)}\bigr)\in\mathbb{R}^n,
\]
where $(i_1,\ldots,i_n)$ is a permutation of $(1,\ldots,n)$; hence,
$T$ is sufficient for $P\in\mathcal{P}$.\par
\vspace{0.5em}
\textbf{Sol.}\par
Let $\mathfrak{S}_n$ denote the set of permutations of
$\{1,\ldots,n\}$, and define
\[
K(t,A)
=\frac{1}{n!}\sum_{\pi\in\mathfrak{S}_n}
\mathbf{1}_A(\pi(t)),
\qquad A\in\mathcal{B}(\R^n).
\]
This is a probability kernel independent of $P$.

For every Borel set $C$ in the state space of $T$, exchangeability
of $X$ and the identity $T(\pi(X))=T(X)$ give
\[
\E[\mathbf{1}_A(X)\mathbf{1}_{\{T\in C\}}]
=
\E[\mathbf{1}_A(\pi(X))\mathbf{1}_{\{T\in C\}}].
\]
Averaging over all permutations yields
\[
\begin{aligned}
\E[\mathbf{1}_A(X)\mathbf{1}_{\{T\in C\}}]
&=
\E\left[
\frac{1}{n!}\sum_{\pi\in\mathfrak{S}_n}
\mathbf{1}_A(\pi(X))\mathbf{1}_{\{T\in C\}}
\right]\\
&=
\E[K(T,A)\mathbf{1}_{\{T\in C\}}],
\end{aligned}
\]
because $X$ and $T$ have the same collection of permuted vectors.
Hence $K(T,A)$ is a version of $\P(X\in A\mid T)$.

Since the coordinates of $T$ are distinct almost surely,
$K(T,\cdot)$ assigns mass $1/n!$ to each of the $n!$
permutations of $T$. This conditional distribution does not
depend on $P$, so $T$ is sufficient.

\subsubsection*{Ex. 2.25}
Show that if $T$ is a sufficient statistic and
\[
T=\psi(S),
\]
where $\psi$ is measurable and $S$ is another statistic, then $S$ is
sufficient.\par
\vspace{0.5em}
\textbf{Sol.}\par
Since $T=\psi(S)$ and $\psi$ is measurable,
\[
\sigma(T)\subseteq\sigma(S).
\]
By sufficiency of $T$, the conditional distribution of $X$ given
$T=t$ has a version $K(t,\mathrm{d}x)$ that does not depend on
the underlying distribution $P$.

Write $S=s(X)$. Under the probability measure
$K(t,\mathrm{d}x)$, let $L(t,u,\mathrm{d}x)$ be a regular
conditional distribution of $X$ given $s(X)=u$.
Since $K$ and $s$ do not depend on $P$, this conditional kernel
can also be chosen independently of $P$.
By conditioning first on $T$ and then on $S$,
\[
L(T,S,\mathrm{d}x)
\]
is a version of the conditional distribution of $X$ given
$(T,S)$.

Because $T=\psi(S)$, conditioning on $(T,S)$ is equivalent
to conditioning on $S$. Hence
\[
L(\psi(S),S,\mathrm{d}x)
\]
is a version of the conditional distribution of $X$ given $S$,
and it does not depend on $P$. Therefore, $S$ is sufficient.



\subsection*{Problem 2. Prediction}
Give an example in which $Z$ determines $Y$ exactly, but
\[
\operatorname{Var}(Z\mid Y)=\operatorname{Var}(Z)
\qquad\text{almost surely}.
\]
Require $Y$ to be nonconstant and $Z$ to have finite, positive variance.
Verify both properties.\par
\vspace{1em}
\textbf{Sol.}\par
Let $Z$ have the distribution
\[
\P(Z=-1)=\P(Z=1)=\frac14,\qquad
\P(Z=-\sqrt{2})=\P(Z=\sqrt{2})=\frac18,\qquad
\P(Z=0)=\frac14.
\]
Define
\[
Y=\mathbf{1}_{\{|Z|=1\}}.
\]
Then $Y$ is a measurable function of $Z$, so $Z$ determines $Y$
exactly. Moreover,
\[
\P(Y=0)=\P(Y=1)=\frac12,
\]
so $Y$ is nonconstant.

By symmetry, $\E Z=0$, and
\[
\operatorname{Var}(Z)=\E Z^2
=\frac14+\frac14
+2\left(\frac18+\frac18\right)=1.
\]
Thus $Z$ has finite, positive variance.

Given $Y=1$, $Z$ takes the values $-1$ and $1$ with equal
probability. Hence
\[
\E[Z\mid Y=1]=0,\qquad
\E[Z^2\mid Y=1]=1.
\]
Given $Y=0$, the conditional probabilities are
\[
\P(Z=-\sqrt{2}\mid Y=0)=\frac14,\qquad
\P(Z=0\mid Y=0)=\frac12,\qquad
\P(Z=\sqrt{2}\mid Y=0)=\frac14.
\]
Therefore,
\[
\E[Z\mid Y=0]=0,\qquad
\E[Z^2\mid Y=0]
=2\cdot\frac14+2\cdot\frac14=1.
\]
Consequently,
\[
\operatorname{Var}(Z\mid Y)
=\E[Z^2\mid Y]-\bigl(\E[Z\mid Y]\bigr)^2
=1
=\operatorname{Var}(Z)
\qquad\text{a.s.}
\]

\subsection*{Problem 3. Exponential families}
Let $X\in\R^d$ have a density, with respect to Lebesgue measure, of the form
\[
f_\theta(x)
=h(x)\exp\!\left\{\theta^\top T(x)-\xi(\theta)\right\},
\qquad \theta\in\Theta.
\]
Here $h$ is nonnegative and Borel measurable,
$T:\R^d\to\R^k$ is continuous, and the natural parameter space is
\[
\Theta
=
\left\{
\theta\in\R^k:
0<\int_{\R^d}h(x)e^{\theta^\top T(x)}\,dx<\infty
\right\}.
\]
Fix $\theta_0$ in its interior, and let $B$ be an open convex neighborhood
of $\theta_0$ contained in $\Theta$.

\begin{enumerate}[label=(\arabic*),leftmargin=*]
\item Show that
\[
\xi(\theta)
=\log\int_{\R^d}h(x)e^{\theta^\top T(x)}\,dx
\]
and that $\xi\in C^\infty(B)$. Thus, for every multi-index
$\alpha=(\alpha_1,\ldots,\alpha_k)\in\{0,1,2,\ldots\}^k$, with
$|\alpha|=\alpha_1+\cdots+\alpha_k$,
\[
\frac{\partial^{|\alpha|}\xi}
{\partial\theta_1^{\alpha_1}\cdots\partial\theta_k^{\alpha_k}}
\]
exists and is continuous on $B$.

\item For $\theta\in B$, express $\E_\theta T(X)$,
$\operatorname{Cov}_\theta(T(X))$, and the moment-generating function of
$T(X)$ using $\xi$ and its derivatives. Specify where the
moment-generating function is finite.

\item Show that $\Theta$ is convex and that $\xi$ is convex on $\Theta$.
Show that $\xi$ is strictly convex on $B$ if $T(X)$ is not supported on any
proper affine hyperplane: there are no $a\ne0$ and $b$ with
\[
a^\top T(X)=b
\qquad\text{almost surely under }\P_{\theta_0}.
\]
\end{enumerate}
\textbf{Sol.}\par
Write
\[
Z(\theta)
:=
\int_{\R^d}h(x)e^{\theta^\top T(x)}\,dx.
\]

\begin{enumerate}[label=(\arabic*),leftmargin=*]
\item Since $f_\theta$ is a probability density,
\[
1
=
\int_{\R^d}f_\theta(x)\,dx
=
e^{-\xi(\theta)}Z(\theta).
\]
Consequently,
\[
\xi(\theta)=\log Z(\theta).
\]

We next justify differentiation under the integral sign.
Fix $\theta^\ast\in B$. Since $B$ is open, choose $a>0$ such that
\[
\theta^\ast+[-2a,2a]^k\subseteq B.
\]
For a multi-index $\alpha$, put
\[
T(x)^\alpha=\prod_{j=1}^k T_j(x)^{\alpha_j}.
\]
There is a constant $C_{\alpha,a}<\infty$ such that, for every
$u\in\R^k$,
\[
|u^\alpha|
\leq C_{\alpha,a}e^{a\sum_{j=1}^k|u_j|}.
\]
Thus, whenever $\theta\in\theta^\ast+[-a,a]^k$,
\[
\begin{aligned}
h(x)|T(x)^\alpha|e^{\theta^\top T(x)}
&\leq
C_{\alpha,a}h(x)
e^{(\theta^\ast)^\top T(x)+2a\sum_{j=1}^k|T_j(x)|}\\
&\leq
C_{\alpha,a}
\sum_{v\in\{-1,1\}^k}
h(x)e^{(\theta^\ast+2av)^\top T(x)}.
\end{aligned}
\]
The right-hand side is integrable, because all the parameters
$\theta^\ast+2av$ belong to $\Theta$.

Differentiation under the integral sign and dominated convergence
therefore give
\[
\partial^\alpha Z(\theta)
=
\int_{\R^d}
T(x)^\alpha h(x)e^{\theta^\top T(x)}\,dx,
\]
and these derivatives are continuous in a neighborhood of
$\theta^\ast$. Since $\theta^\ast$ was arbitrary,
$Z\in C^\infty(B)$. As $Z>0$ on $B$, it follows that
\[
{\xi=\log Z\in C^\infty(B)}.
\]

\item Differentiating $\xi=\log Z$, we obtain
\[
\frac{\partial\xi}{\partial\theta_j}(\theta)
=
\frac{
\int T_j(x)h(x)e^{\theta^\top T(x)}\,dx
}{
Z(\theta)
}
=
\E_\theta T_j(X).
\]
Hence
\[
{\E_\theta T(X)=\nabla\xi(\theta)}.
\]

Similarly,
\[
\begin{aligned}
\frac{\partial^2\xi}
{\partial\theta_i\partial\theta_j}(\theta)
&=
\frac{
\int T_i(x)T_j(x)h(x)e^{\theta^\top T(x)}\,dx
}{
Z(\theta)
}\\
&\quad-
\frac{
\int T_i(x)h(x)e^{\theta^\top T(x)}\,dx
}{
Z(\theta)
}
\frac{
\int T_j(x)h(x)e^{\theta^\top T(x)}\,dx
}{
Z(\theta)
}\\
&=
\operatorname{Cov}_\theta(T_i(X),T_j(X)).
\end{aligned}
\]
Therefore,
\[
{
\operatorname{Cov}_\theta(T(X))
=
\nabla^2\xi(\theta)
}.
\]

For $u\in\R^k$, the moment-generating function of $T(X)$ is
\[
\begin{aligned}
M_\theta(u)
&:=\E_\theta e^{u^\top T(X)}\\
&=
e^{-\xi(\theta)}
\int_{\R^d}h(x)e^{(\theta+u)^\top T(x)}\,dx.
\end{aligned}
\]
It is finite exactly when $\theta+u\in\Theta$, and in that case
\[
{
M_\theta(u)
=
\exp\{\xi(\theta+u)-\xi(\theta)\}
}.
\]
In particular, it is finite in a neighborhood of $u=0$ whenever
$\theta\in B$.

\item Let $\theta_1,\theta_2\in\Theta$ and $s\in(0,1)$.
By H\"older's inequality,
\[
\begin{aligned}
Z((1-s)\theta_1+s\theta_2)
&=
\int_{\R^d}
\bigl(h(x)e^{\theta_1^\top T(x)}\bigr)^{1-s}
\bigl(h(x)e^{\theta_2^\top T(x)}\bigr)^s\,dx\\
&\leq
Z(\theta_1)^{1-s}Z(\theta_2)^s
<\infty.
\end{aligned}
\]
This integral is also positive: $h$ is positive on a set of
positive Lebesgue measure, and the exponential factor is
strictly positive. Hence
\[
(1-s)\theta_1+s\theta_2\in\Theta,
\]
so $\Theta$ is convex.

Taking logarithms gives
\[
\xi((1-s)\theta_1+s\theta_2)
\leq
(1-s)\xi(\theta_1)+s\xi(\theta_2).
\]
Therefore, $\xi$ is convex on $\Theta$.

To prove strict convexity on $B$, observe that all the laws
$\P_\theta$, $\theta\in\Theta$, are mutually absolutely continuous:
their densities are strictly positive on the same set
$\{x:h(x)>0\}$.

For every $\theta\in B$ and every nonzero $a\in\R^k$,
part (2) gives
\[
a^\top\nabla^2\xi(\theta)a
=
\operatorname{Var}_\theta(a^\top T(X)).
\]
If this variance were zero, there would be a constant $b$ such that
\[
a^\top T(X)=b
\qquad \P_\theta\text{-a.s.}
\]
Mutual absolute continuity would imply the same identity
$\P_{\theta_0}$-a.s., contradicting the hypothesis.
Thus
\[
a^\top\nabla^2\xi(\theta)a>0
\qquad
\text{for every }\theta\in B,\ a\ne0.
\]

Finally, for distinct $\theta_1,\theta_2\in B$, define
\[
g(s)=\xi(\theta_1+s(\theta_2-\theta_1)),
\qquad 0\leq s\leq1.
\]
Since $B$ is convex,
\[
g''(s)
=
(\theta_2-\theta_1)^\top
\nabla^2\xi(\theta_1+s(\theta_2-\theta_1))
(\theta_2-\theta_1)
>0.
\]
Consequently, for $0<s<1$,
\[
\xi((1-s)\theta_1+s\theta_2)
<
(1-s)\xi(\theta_1)+s\xi(\theta_2).
\]
Hence $\xi$ is strictly convex on $B$.
\end{enumerate}

\vspace{1em}

\subsection*{Problem 4. Bias--variance tradeoff}
Let
\[
X_1,\ldots,X_n\overset{\mathrm{i.i.d.}}{\sim}\operatorname{Bernoulli}(\theta),
\qquad \theta\in[0,1],
\]
and write
\[
\overline X=\frac{1}{n}\sum_{i=1}^n X_i.
\]
Consider the estimator
\[
\widehat\theta
=(1-\varepsilon_n)\left(\overline X-\frac12\right)+\frac12,
\qquad 0<\varepsilon_n<1.
\]
The constant $\varepsilon_n$ may depend only on $n$. Setting
$\varepsilon_n=0$ would recover the sample mean.

\begin{enumerate}[label=(\arabic*),leftmargin=*]
\item Find the bias and variance of $\widehat\theta$.

\item Find the choice of $\varepsilon_n\in(0,1)$ that minimizes the
worst-case mean squared error
\[
\sup_{\theta\in[0,1]}
\E_\theta\!\left[(\widehat\theta-\theta)^2\right].
\]
\end{enumerate}
\vspace{1em}
\textbf{Sol.}\par
\begin{enumerate}[label=(\arabic*),leftmargin=*]
\item Since $X_1,\ldots,X_n$ are i.i.d.\ Bernoulli$(\theta)$,
\[
\E_\theta\overline X=\theta,
\qquad
\operatorname{Var}_\theta(\overline X)
=\frac{\theta(1-\theta)}{n}.
\]
Therefore,
\[
\E_\theta\widehat\theta
=(1-\epsilon_n)\left(\theta-\frac12\right)+\frac12
=(1-\epsilon_n)\theta+\frac{\epsilon_n}{2}.
\]
Hence the bias is
\[
b_\theta(\widehat\theta)
=\E_\theta\widehat\theta-\theta
=\epsilon_n\left(\frac12-\theta\right).
\]
The variance is
\[
\operatorname{Var}_\theta(\widehat\theta)
=(1-\epsilon_n)^2
\operatorname{Var}_\theta(\overline X)
=\frac{(1-\epsilon_n)^2}{n}\theta(1-\theta).
\]

\item By the bias--variance decomposition,
\[
\E_\theta[(\widehat\theta-\theta)^2]
=
\frac{(1-\epsilon_n)^2}{n}\theta(1-\theta)
+\epsilon_n^2\left(\theta-\frac12\right)^2.
\]
Using
\[
\theta(1-\theta)
=\frac14-\left(\theta-\frac12\right)^2,
\]
we obtain
\[
\E_\theta[(\widehat\theta-\theta)^2]
=
\frac{(1-\epsilon_n)^2}{4n}
+
\left(
\epsilon_n^2-\frac{(1-\epsilon_n)^2}{n}
\right)
\left(\theta-\frac12\right)^2.
\]
Since $(\theta-\frac12)^2$ ranges over $[0,1/4]$,
\[
\phi(\epsilon_n)
:=
\sup_{\theta\in[0,1]}
\E_\theta[(\widehat\theta-\theta)^2]
=
\frac14\max\left\{
\epsilon_n^2,\frac{(1-\epsilon_n)^2}{n}
\right\}.
\]
Since $0<\epsilon_n<1$, this gives
\[
\sqrt n\,\epsilon_n=1-\epsilon_n,
\qquad
\epsilon_n^\ast=\frac{1}{1+\sqrt n}.
\]
The resulting worst-case mean squared error is
\[
\phi(\epsilon_n^\ast)
=\frac{1}{4(1+\sqrt n)^2}
.
\]
\end{enumerate}
\end{document}
